\chapter{График движения}
\label{ch:schedule}

В таблице~\ref{tab:schedule} приведён график движения по дням. Хронометраж чисто ходового времени (ЧХВ) указан в часах и минутах. Протяжённость участков дана по карте (с учётом коэффициента 1,2 для горной местности).

\begin{schedulestable}[tab:schedule]{График движения}
16.07 & 1 & Пос. Турасу — дол. р. Тон          & 8,5 & 3:15 & +450  & –                                   & Ясно, +25° \\ \hline
17.07 & 2 & дол. р. Тон — дол. р. Зындан       & 12,2& 5:00 & +620  & –                                   & Переменная облачность \\ \hline
18.07 & 3 & дол. р. Зындан — пер. 267           & 9,0 & 6:30 & +870  & \makecell{пер. 267 (1Б, 4086),\\осыпь 20-25°} & Облачно, ветер \\ \hline
19.07 & 4 & пер. 267 — пер. Кельторташ          & 10,5& 7:15 & +750/-400 & пер. Кельторта (1Б, 4117), снежный склон & Снегопад \\ \hline
% ... остальные строки
\end{schedulestable}

*Протяжённость указана с учётом коэффициента 1,2 для горной местности.